\documentclass[11pt]{article}
\usepackage[utf8]{inputenc}
\usepackage{amsmath, amssymb, amsthm}
\usepackage{geometry}
\geometry{margin=1in}

\newtheorem{theorem}{Theorem}
\newtheorem{definition}{Definition}
\newtheorem{lemma}{Lemma}

\title{Change of Variables in Probability}
\author{Math Foundations \textendash{} Concept 28}
\date{}

\begin{document}
\maketitle

\subsection*{First, an example}

Before the general theory, here is the prototype. Let $X \sim \mathrm{Uniform}(0,1)$
and define $Y = -\log X$. The transformation $g(x) = -\log x$ has inverse
$g^{-1}(y) = e^{-y}$ and derivative magnitude $|g'(x)| = 1/x$. Plugging into the
density-transformation rule (Theorem~\ref{thm:1d} below) with $f_X \equiv 1$ on $(0,1)$,
\[
f_Y(y) \;=\; f_X(e^{-y}) \cdot e^{-y} \;=\; e^{-y}, \qquad y \ge 0.
\]
\emph{Read aloud:} ``f sub Y of y equals f sub X of e to the minus y, times e to the minus y, which equals e to the minus y, for y greater than or equal to zero.''
So $Y$ is exponential with rate $1$. The factor $e^{-y}$ is the Jacobian of the inverse
map; everything that follows is a careful generalization of this single calculation.

\section{Setup}

Let $X$ be a continuous random variable taking values in an open set $U \subseteq \mathbb{R}^n$
with probability density function (PDF) $f_X : U \to [0,\infty)$. Let $g : U \to V$
be a transformation; we wish to find the density of $Y = g(X)$.

\begin{definition}[Diffeomorphism]
A map $g : U \to V$ between open subsets of $\mathbb{R}^n$ is a $C^1$ \emph{diffeomorphism}
if $g$ is a bijection, $g \in C^1(U)$, and $g^{-1} \in C^1(V)$. Equivalently, $g$ is $C^1$
and its Jacobian matrix $J_g(x) = (\partial g_i/\partial x_j)_{ij}$ is invertible for every $x \in U$.
\end{definition}

\begin{definition}[Push-forward density]
If $X$ has density $f_X$ and $Y = g(X)$ has a density, we call that density the
\emph{push-forward} of $f_X$ along $g$ and denote it $f_Y = g_\# f_X$.
\end{definition}

\section{Main results}

\begin{theorem}[1D change of variables]
\label{thm:1d}
Let $X$ be a continuous random variable with density $f_X$ on an open interval
$(a,b) \subseteq \mathbb{R}$. Let $g : (a,b) \to (c,d)$ be a $C^1$ function with
$g'(x) \neq 0$ for all $x \in (a,b)$. Then $g$ is strictly monotone, hence bijective,
and $Y = g(X)$ has density
\[
f_Y(y) \;=\; \frac{f_X\!\bigl(g^{-1}(y)\bigr)}{\bigl|g'\!\bigl(g^{-1}(y)\bigr)\bigr|},
\qquad y \in (c,d).
\]
\emph{Read aloud:} ``f sub Y of y equals f sub X of g-inverse of y, divided by the absolute value of g-prime of g-inverse of y, for y in the open interval c to d.''
\end{theorem}

\begin{proof}
Because $g'$ is continuous and never zero on the connected interval $(a,b)$, either
$g' > 0$ everywhere or $g' < 0$ everywhere. We treat the two monotone cases via the
cumulative distribution function (CDF) approach.

\medskip
\noindent\textbf{Case 1: $g$ strictly increasing.}
Then $g : (a,b) \to (c,d)$ is a bijection with $C^1$ inverse $g^{-1} : (c,d) \to (a,b)$,
and for any $y \in (c,d)$,
\[
F_Y(y) \;=\; P(Y \le y) \;=\; P\bigl(g(X) \le y\bigr) \;=\; P\bigl(X \le g^{-1}(y)\bigr) \;=\; F_X\!\bigl(g^{-1}(y)\bigr).
\]
\emph{Read aloud:} ``capital F sub Y of y equals the probability that Y is at most y, equals the probability that g of X is at most y, equals the probability that X is at most g-inverse of y, equals capital F sub X of g-inverse of y.''
Differentiating in $y$ (chain rule; $F_X$ is differentiable with $F_X' = f_X$, and
$g^{-1}$ is differentiable with $(g^{-1})'(y) = 1/g'(g^{-1}(y))$ by the inverse function
theorem),
\[
f_Y(y) \;=\; F_Y'(y) \;=\; f_X\!\bigl(g^{-1}(y)\bigr)\,\bigl(g^{-1}\bigr)'(y)
\;=\; \frac{f_X\!\bigl(g^{-1}(y)\bigr)}{g'\!\bigl(g^{-1}(y)\bigr)}.
\]
\emph{Read aloud:} ``f sub Y of y equals capital F sub Y prime of y, equals f sub X of g-inverse of y, times the derivative of g-inverse at y, equals f sub X of g-inverse of y divided by g-prime of g-inverse of y.''
Since $g' > 0$, we have $|g'| = g'$, so this equals the claimed formula.

\medskip
\noindent\textbf{Case 2: $g$ strictly decreasing.}
Then for $y \in (c,d)$,
\[
F_Y(y) \;=\; P\bigl(g(X) \le y\bigr) \;=\; P\bigl(X \ge g^{-1}(y)\bigr) \;=\; 1 - F_X\!\bigl(g^{-1}(y)\bigr).
\]
\emph{Read aloud:} ``capital F sub Y of y equals the probability that g of X is at most y, equals the probability that X is at least g-inverse of y, equals one minus capital F sub X of g-inverse of y.''
Differentiating,
\[
f_Y(y) \;=\; -\, f_X\!\bigl(g^{-1}(y)\bigr)\,\bigl(g^{-1}\bigr)'(y)
\;=\; -\, \frac{f_X\!\bigl(g^{-1}(y)\bigr)}{g'\!\bigl(g^{-1}(y)\bigr)}.
\]
\emph{Read aloud:} ``f sub Y of y equals negative f sub X of g-inverse of y times the derivative of g-inverse at y, equals negative f sub X of g-inverse of y divided by g-prime of g-inverse of y.''
Since $g' < 0$, $-1/g' = 1/|g'|$, recovering the claimed formula.

\medskip
In both cases,
$f_Y(y) = f_X(g^{-1}(y)) / |g'(g^{-1}(y))|$, completing the proof. \qedhere
\end{proof}

\begin{theorem}[Multivariate change of variables]
\label{thm:multi}
Let $X$ be a continuous random vector on an open set $U \subseteq \mathbb{R}^n$ with
density $f_X$, and let $g : U \to V$ be a $C^1$ diffeomorphism onto an open set
$V \subseteq \mathbb{R}^n$. Then $Y = g(X)$ has density
\[
f_Y(y) \;=\; f_X\!\bigl(g^{-1}(y)\bigr) \cdot \bigl|\det J_{g^{-1}}(y)\bigr|
\;=\; \frac{f_X(x)}{\bigl|\det J_g(x)\bigr|}\bigg|_{x = g^{-1}(y)},
\qquad y \in V.
\]
\emph{Read aloud:} ``f sub Y of y equals f sub X of g-inverse of y, times the absolute value of the determinant of the Jacobian of g-inverse at y, which also equals f sub X of x divided by the absolute value of the determinant of the Jacobian of g at x, evaluated at x equals g-inverse of y, for y in V.''
\end{theorem}

\begin{proof}[Proof sketch]
For any Borel set $B \subseteq V$, using $g^{-1}(B) \subseteq U$ and the multivariate
substitution theorem from analysis,
\[
P(Y \in B) = P(X \in g^{-1}(B)) = \int_{g^{-1}(B)} f_X(x)\, dx
= \int_B f_X(g^{-1}(y))\,|\det J_{g^{-1}}(y)|\,dy.
\]
\emph{Read aloud:} ``the probability that Y lies in B equals the probability that X lies in g-inverse of B, equals the integral over g-inverse of B of f sub X of x, d-x, equals the integral over B of f sub X of g-inverse of y, times the absolute value of the determinant of the Jacobian of g-inverse at y, d-y.''
Because $B$ was arbitrary, the integrand is the density of $Y$.
\end{proof}

\section{Canonical example}

Let $X \sim \mathrm{Uniform}[0,1]$ and $Y = -\log X$. Then $g(x) = -\log x$ is
$C^1$ on $(0,1)$ with $g'(x) = -1/x$, $|g'(x)| = 1/x$, and $g^{-1}(y) = e^{-y}$.
By Theorem~\ref{thm:1d},
\[
f_Y(y) \;=\; \frac{1}{1/e^{-y}} \;=\; e^{-y}, \qquad y > 0,
\]
\emph{Read aloud:} ``f sub Y of y equals one divided by one over e to the minus y, which equals e to the minus y, for y greater than zero.''
i.e.\ $Y \sim \mathrm{Exp}(1)$. This identity underlies inverse-CDF sampling.

\end{document}
